%! Expressions and Operations — Unit 1, Lesson 1.1 — Warm-Up
\documentclass[10pt]{article}
\usepackage{algebra2-article}
\usepackage{algebra2-boxes}
\newcommand{\TallMath}[1]{$\displaystyle #1\rule[-1.4em]{0pt}{3.2em}$}

\begin{document}

\pageheader{Unit 1, Lesson 1.1}{Warm-Up}

\vspace{0.08in}

\begin{notesbox}{Getting Ready: Skills We'll Use Today}
\small These three skills power today's review of expressions. Work quickly.

\vspace{4pt}
\textbf{1. Order of operations.} Evaluate each one.
\begin{center}
\renewcommand{\arraystretch}{1.5}
\begin{tabular}{@{}l l l@{}}
(a) \; $12-3(5-7)=$ \blank{1.8cm} &
(b) \; $|{-6}|+\sqrt{16}-\sqrt[3]{27}=$ \blank{1.8cm} &
(c) \; $\dfrac{4+2^3}{6}=$ \blank{1.8cm} \\
\end{tabular}
\end{center}
Now the one that always starts an argument. Let $x=-3$.
\begin{center}
\renewcommand{\arraystretch}{1.5}
\begin{tabular}{@{}l l@{}}
(d) \; $x^2=$ \blank{1.6cm} &
(e) \; $-x^2=$ \blank{1.6cm} \\
\end{tabular}
\end{center}
Those two are \emph{not} the same. What exactly is the exponent attached to in each one? \writeline

\vspace{4pt}
\textbf{2. Exponents --- write them out and count.} Do \emph{not} use a rule. Expand first.
\begin{center}
\renewcommand{\arraystretch}{1.6}
\begin{tabularx}{\linewidth}{@{}l X l@{}}
\hline
\textbf{Expression} & \textbf{Written out as factors} & \textbf{Single power} \\
\hline
$2^5\cdot 2^3$ & \blank{7.0cm} & \blank{1.6cm} \\
$\dfrac{a^6}{a^2}$ & \blank{7.0cm} & \blank{1.6cm} \\
$(a^3)^2$ & \blank{7.0cm} & \blank{1.6cm} \\
\hline
\end{tabularx}
\end{center}
Look at your three answers. Where did the exponents' $+$, $-$, and $\times$ \emph{come from}?
\writeline

\vspace{4pt}
\textbf{3. Distributing --- twice.}
\begin{center}
\begin{minipage}[t]{0.44\linewidth}
\centering (a) \; $4(2x-5)=$ \blank{2.6cm}
\vspace{0.35cm}
\end{minipage}
\begin{minipage}[t]{0.52\linewidth}
\centering
(b) \; Fill in the area model for $(x+3)(x+2)$.

\vspace{3pt}
\renewcommand{\arraystretch}{1.6}
\begin{tabular}{|c|c|c|}
\hline
$\times$ & $x$ & $+3$ \\ \hline
$x$ & \blank{1.2cm} & \blank{1.2cm} \\ \hline
$+2$ & \blank{1.2cm} & \blank{1.2cm} \\ \hline
\end{tabular}

\vspace{3pt}
Total: \blank{3.4cm}
\end{minipage}
\end{center}
Both parts used the \emph{same} property from Lesson 1.0. Name it. \blank{5.0cm}

\end{notesbox}

\end{document}
